\section{Related Works}
\label{sec:related}

% \todo{refine and prune}

\subsection{Agentic RL}
% \todo{tool orchestra, agent lightning, arpo, etc etc}
Agentic language models interleave natural language with grounded actions such as tool calls, code execution, web navigation, and file edits, bootstrapped with supervised trajectories and prompting scaffolds~\citep{yao2023react,schick2023toolformer}.
Recent work also investigates explicit \emph{tool/model orchestration}, where a lightweight controller selects among tools or specialist models to improve efficiency and capability \citep{su2025toolorchestra}.
On the data and evaluation side, API-interaction corpora such as ToolBench and ToolLLM \citep{xu2023toolbench,qin2023toolllm} are complemented by benchmarks spanning embodied/text-world environments~\citep{shridhar2020alfworld,wang2022scienceworld}, web interaction~\citep{yao2022webshop}, multi-domain agent evaluation suites~\citep{ma2024agentboard}, and social-agent environments~\citep{wang2024sotopiapi}.
In software engineering, execution-grounded agents are trained and assessed on real repositories and environment-centric setups~\citep{wang2025openhandsopenplatformai,pan2025trainingsoftwareengineeringagents}.
Simulated multi-turn data generation has become a common ingredient for scaling agentic supervision \citep{prabhakar2025apigenmtagenticpipelinemultiturn}, and recent theory work begins to formalize why tool-integrated reasoning can expand capability \citep{lin2025tir}.

Reinforcement learning optimizes an interaction policy against task-derived or preference-derived feedback~\citep{ouyang2022training,schulman2017ppo,schulman2015trpo}.
For long-horizon agents, recent work targets exploration and credit assignment via stepwise progress attribution, exploration-based trajectory optimization, and verifier/meta-reasoning rewards \citep{wang2025sparl,song2024trialerror,zhang2025rlvmr}, as well as tree-structured sampling/search and robustness to imperfect process verifiers \citep{ji2025treesearchllmrl,rohatgi2025tamingprocessverifiers}.
End-to-end multi-turn RL has been applied directly to tool-integrated reasoning and multi-tool use \citep{xue2025simpletir,dong2025toolstar}, with systems advances for scalable and stable policy optimization \citep{yang2025dcpo,yu2025dapo}.
Large-scale agent training and environment design further improve the efficiency and stability of multi-turn RL for tool use \citep{luo2025agentlightningtrainai,dong2025agenticreinforcedpolicyoptimization}.
Multi-agent extensions study coordination and credit assignment in interactive settings \citep{zhou2025sweetrl,xi2025agentgymrl}.
In parallel, RL with \emph{verifiable} rewards (RLVR) has driven large gains in domains with automatic checkers, such as math and theorem proving \citep{shao2024deepseekmath,xin2024deepseekproverv15harnessingproofassistant,Guo_2025}.
Our work sits at this intersection: we target agentic interaction (terminal, code, browsing) while leveraging \emph{verifiability} by converting supervised trajectories into \emph{turn-level} rewards via completion-based verifiers, enabling on-policy optimization without training a separate reward model.

\subsection{Process Guidance}
% \todo{ArCHer (from levine's lab), SweetRL (from Levine's Lab), on policy distillation, Reinforcement Learning via Self-Distillation, GLM 4.5 (formatting + a), Process Reward Models, ARPO, DeepSeek process guidance}

A key challenge in long-horizon reasoning and acting is credit assignment: many failures are attributable to intermediate decisions rather than the final output.
Inference-time guidance encourages explicit intermediate reasoning, such as chain-of-thought prompting \citep{wei2022cot}, self-consistency decoding \citep{wang2022selfconsistency}, and tree-structured deliberation \citep{yao2023treeofthoughts}.
For agents, ReAct can be viewed as a lightweight process scaffold that interleaves intermediate rationales with tool actions \citep{yao2023react}.

Training-time approaches seek denser supervision signals.
Process reward models (PRMs) and step-level verification score intermediate reasoning trajectories rather than only final answers \citep{lightman2023letsverifystepstep}, and automated process-supervision pipelines aim to synthesize intermediate labels or critiques at scale \citep{luo2024improvemathematicalreasoninglanguage}.
In agentic RL, several lines of work obtain denser credit assignment through verifier and search-guided optimization~\citep{dong2025agenticreinforcedpolicyoptimization,ji2025treesearchllmrl}, and through explicit modeling of imperfect process verifiers \citep{rohatgi2025tamingprocessverifiers}.
Trajectory densification and multi-turn RL systems such as Agent Lightning provide another practical route to stabilize exploration and learning in tool-use environments \citep{luo2025agentlightningtrainai}.
Hierarchical multi-turn RL frameworks such as ArCHer~\citep{zhou2024archer} explicitly separate turn-level and token-level learning signals to improve multi-turn credit assignment, while social-agent RL methods such as SweetRL~\citep{zhou2025sweetrl} study protocol-level learning and credit assignment in multi-agent interaction.

Complementary to explicit reward shaping, distillation-based post-training uses on-policy rollouts to produce more robust or efficient policies \citep{lu2025onpolicydistillation}.
Recent work on reinforcement learning via self-distillation converts rich textual feedback (e.g., error messages) into dense learning signals without external teachers \citep{hubotter2026selfdistillation}.
Model development reports for agentic/coding systems (e.g., GLM-4.5) emphasize strict action formats and multi-stage post-training to stabilize long-horizon behavior \citep{5team2025glm45agenticreasoningcoding}, and large-scale RLVR systems such as DeepSeek-R1 highlight the effectiveness of enforcing verifiable outcomes for complex reasoning \citep{Guo_2025}.
PivotRL contributes a complementary perspective for agentic tasks with demonstrations: instead of learning a separate PRM, we treat each tool boundary (pivot) as a decision point and use the demonstration completion at that pivot as a \emph{verifier-derived} dense reward, which functions as process guidance tailored to multi-turn interaction.

\subsection{The Role of SFT in Agentic Post-Training}
% \todo{superficial alignment hypothesis: https://arxiv.org/pdf/2602.15829}

The role of supervised fine-tuning (SFT) in post-training has been actively debated.
\emph{Superficial alignment hypothesis} posits that alignment  merely surfaces the knowledge learned almost entirely during pre-training~\citep{zhou2023limaalignment}. Recent work operationalizes this perspective via explicit task-complexity controls \citep{vergarabrowne2026sahcomplexity} and revisits its empirical predictions across settings \citep{raghavendra2024revisitingSAH}.
These observations motivate training on  carefully curated trajectories, which is adopted and further advanced by PivotRL.

As a major class of alternatives, \emph{preference optimization} optimizes against preference comparisons.
Direct Preference Optimization (DPO) reparameterizes RLHF objectives into a supervised-style loss \citep{rafailov2023dpo}, and subsequent work explores variants that modify the implicit reward, regularization, or reference-model assumptions \citep{ethayarajh2024kto,hong2024orpo,meng2024simpo}, including segment-level formulations for multi-turn social agents \citep{kong2025sdpo}.
In contrast, PivotRL keeps demonstration data but eliminates \emph{behavior cloning as the training objective} for agentic trajectories: we use demonstrations only to define \emph{verifiable turn-level rewards} and then optimize on-policy.
This positioning is aligned with evidence that RL-style objectives can improve generalization relative to pure imitation when rewards reshape or correct the supervision signal \citep{chu2025sftmemorizesrlgeneralizes,wu2025generalization}.


\subsection{Adapting Behavior Cloning into RL}
% \todo{Dagger maybe}

% \todo{Prefix-RL (concurrent work)}

PivotRL can be viewed through the lens of imitation learning and RL-from-demonstrations.
Classical behavior cloning can suffer from compounding errors under distribution shift, motivating on-policy data aggregation methods such as DAgger \citep{ross2011dagger}, occupancy-measure matching approaches such as GAIL \citep{ho2016gail}, and RL-from-demonstrations methods that combine expert data with RL updates such as DQfD \citep{hester2018dqfd}.
Recent work also explores using partial successful traces as \emph{prefixes} to guide on-policy optimization; for example, PrefixRL conditions RL on off-policy prefixes to improve exploration on hard problems \citep{setlur2026prefixrl}.

Our approach is closest in spirit to this imitation-to-RL bridge, but tailored to multi-turn tool use.
We start from supervised trajectories, recast each turn as a decision point with a verifiable target action (the demonstration completion), and apply KL-regularized trust-region policy optimization \citep{schulman2017ppo,schulman2015trpo}.
Unlike offline preference optimization methods \citep{rafailov2023dpo,meng2024simpo}, we collect on-policy rollouts and use verifier rewards for turn-level credit assignment; unlike learned reward-model RLHF \citep{ouyang2022training,bai2022helpfulharmless}, we do not require training a separate reward model for these agentic tasks.



\subsection{Mitigating catastrophic forgetting in SFT}
% \todo{RL’S RAZOR: WHY ONLINE REINFORCEMENT
% LEARNING FORGETS LESS, RETAINING BY DOING: THE ROLE OF
% ON-POLICY DATA IN MITIGATING FORGETTING}
% \todo{https://arxiv.org/pdf/2510.14865,https://arxiv.org/pdf/2510.18874,https://arxiv.org/pdf/2405.09673}

A persistent issue in LLM post-training is catastrophic forgetting: optimizing on a narrow target distribution can degrade general capabilities.
Recent empirical studies document forgetting during continual instruction tuning \citep{luo2023empiricalCF}, and analyses link forgetting to optimization properties such as loss-landscape sharpness \citep{li2024revisitingCF}.
Mechanistic work further identifies circuit and component-level contributors to forgetting in transformer fine-tuning \citep{imanov2026mechanisticCF}.
These findings connect to the broader continual-learning literature, which proposes regularization and consolidation strategies such as EWC \citep{kirkpatrick2017ewc} and distillation-based retention methods like Learning without Forgetting \citep{li2016lwf}.

Several practical mitigations are especially relevant for LLMs: parameter-efficient tuning can reduce representational drift and empirically lessen forgetting \citep{biderman2024lora}, while data-mixing ``midtraining'' phases smooth distribution shifts between pretraining and post-training \citep{liu2025midtraining}.
Recent work also argues that online/on-policy RL can forget less than purely offline updates, retaining capabilities through continued interaction \citep{shenfeld2025rlsrazoronlinereinforcement,chen2025retainingdoingroleonpolicy}.
PivotRL complements these approaches from a KL projection perspective: it steers the policy toward task-specific behaviors while remaining close to the initialization and preserving the distribution over task-irrelevant actions, which may help mitigate forgetting when adapting LLM agents to long-horizon tool-use tasks.

% Experiences gained from the verified rewards setting \citep{shao2024deepseekmath,lambert2025tulu3pushingfrontiers,pyatkin2025generalizingverifiableinstructionfollowing}, where the only reward signal is from the final outcome, there is no trivial way to label each action with a process reward. Existing approaches that do not involve human labeling include using a explicit Process Reward Model \citep{lightman2023letsverifystepstep}, Monte Carlo Tree Search \citep{xin2024deepseekproverv15harnessingproofassistant,luo2024improvemathematicalreasoninglanguage}, heuristics-based search \citep{dong2025agenticreinforcedpolicyoptimization} or intermdiate state densification \citep{luo2025agentlightningtrainai} \todo{add more citations}.

% only small amounts of SFT is actually needed for training  \citep{zhou2023limaalignment,chu2025sftmemorizesrlgeneralizes}.